% =============================================================
%  Section: Implementation
% =============================================================

This section describes the engineering decisions that turn the YOLOv1
formulation of Section~\ref{sec:theoretical_framework} into a runnable PyTorch system. We cover, in order,
the data preparation pipeline, the model architecture, the loss
implementation, the training infrastructure and the evaluation pipeline. The
specific bugs that affected V1--V4 are not discussed here; they belong to the
narrative of Section~\ref{sec:bug-investigation}.

% -------------------------------------------------------------
\subsection{Dataset preparation}
\label{sec:impl:dataset}

The training pipeline assumes a single consolidated, cleaned and pre-split
dataset rooted at \texttt{data/coco/}. Constructing it from a freshly
downloaded COCO~2017 archive is a three-step process that we run once,
offline, before the first training run.

\paragraph{Step 1 -- consolidation.} The COCO 2017 archive ships as two
disjoint image folders (\texttt{train2017/} with 118\,287 images and
\texttt{val2017/} with 5\,000 images) and two corresponding annotation files.
\texttt{scripts/consolidate\_dataset.py} merges them into a single flat
\texttt{images/} directory of $123\,287$ images and a single
\texttt{instances\_all.json}. COCO filenames are globally unique, so the move
is collision-free.

\paragraph{Step 2 -- cleaning.} The shipped annotation file is approximately
$550$\,MB; most of that mass is segmentation polygons, captions and licence
metadata that we do not use. \texttt{scripts/clean\_dataset.py} reduces each
record to the minimum the dataset class needs: \texttt{(id, file\_name, width,
height)} per image, \texttt{(image\_id, bbox, category\_id)} per annotation,
and \texttt{(id, name)} per category. In the same pass it removes the
$10\,498$ annotations marked \texttt{iscrowd=1}, which are intended as
``ignore'' regions and would confuse a single-box-per-cell detector. The
cleaned file is approximately seven times smaller and loads under three
seconds on the H200's NVMe.

\paragraph{Step 3 -- deterministic split.}
\texttt{scripts/prepare\_splits.py} performs a single random
$90\,\%/5\,\%/5\,\%$ split with seed~$42$, producing three COCO-format files:
\texttt{instances\_train.json} ($110\,957$ images), \texttt{instances\_val.json}
($6\,165$ images) and \texttt{instances\_test.json} ($6\,165$ images). The
seed is fixed so that the test split is genuinely held out across all six
training runs; any reshuffling between V1 and V6 would have made it impossible
to compare their detection metrics on common terms.

\paragraph{Annotation conversion.} Inside the dataset class
(\texttt{src/dataset.py}), each COCO bounding box
$(x_{\text{tl}}, y_{\text{tl}}, w_{\text{px}}, h_{\text{px}})$ in
\emph{absolute pixel} coordinates is converted to the YOLO grid format used
throughout the rest of the pipeline:
\begin{align}
  x_{\text{norm}} &= \frac{x_{\text{tl}} + w_{\text{px}}/2}{W_{\text{img}}}, &
  y_{\text{norm}} &= \frac{y_{\text{tl}} + h_{\text{px}}/2}{H_{\text{img}}}, \\
  w_{\text{norm}} &= \frac{w_{\text{px}}}{W_{\text{img}}}, &
  h_{\text{norm}} &= \frac{h_{\text{px}}}{H_{\text{img}}}.
\end{align}

\noindent The bounding box is then assigned to the unique grid cell
$(i, j) = (\lfloor S\,y_{\text{norm}} \rfloor, \lfloor S\,x_{\text{norm}} \rfloor)$
whose centre coordinates inside that cell are
$x_{\text{cell}} = S\,x_{\text{norm}} - j$ and
$y_{\text{cell}} = S\,y_{\text{norm}} - i$. The width and height are kept
image-relative; this asymmetry between $(x_{\text{cell}}, y_{\text{cell}})$
and $(w_{\text{norm}}, h_{\text{norm}})$ is the source of Bug~2 discussed in
Section~\ref{sec:bug:wh-decoding}.

\paragraph{One object per cell.} When two ground-truth box centres fall in
the same cell, the dataset retains only the first one encountered in the
iteration order. This implements the architectural constraint stated by
\citet{redmon2016you} that ``each grid cell predicts only one set of class
probabilities, regardless of the number of boxes $B$.'' We do not attempt to
detect or report such collisions during training; on COCO they are
infrequent at $S=7$ relative to the total annotation count, but they account
for a non-trivial fraction of the missed detections discussed in
Section~\ref{sec:results:perclass}.

\paragraph{Input transform.} Because the backbone in V4--V6 is initialised
from torchvision's ImageNet-pretrained ResNet18, the input tensor must follow
the ImageNet normalisation convention. The dataset applies
$\mu_{\text{RGB}} = (0.485,\,0.456,\,0.406)$ and
$\sigma_{\text{RGB}} = (0.229,\,0.224,\,0.225)$ after \texttt{ToTensor}
\citep{krizhevsky2012imagenet,russakovsky2015imagenet}. Visualisation code in
\texttt{src/visualization.py} undoes this normalisation before rendering with
PIL; failing to do so produces tensors with the right shape but visually
incoherent colour distributions, which would otherwise hide localisation
problems in qualitative review.

% -------------------------------------------------------------
\subsection{Model architecture}
\label{sec:impl:model}

The architecture used in V4, V5 and V6 is a three-stage pipeline: a
pretrained ResNet18 convolutional trunk, a learnt adapter that bridges the
backbone's spatial resolution to the detector head, and a fully connected
detection head whose final layer matches the YOLO grid output. The full
data-flow is summarised in Figure~\ref{fig:arch}.

\begin{figure}[!ht]
\centering
\begin{tikzpicture}[
    node distance=4mm,
    every node/.style={font=\small},
    block/.style={draw, rounded corners, align=center, minimum width=7.8cm, minimum height=8mm, inner sep=2pt},
    bb/.style    ={block, fill=blue!8,   draw=blue!50!black},
    adapter/.style={block, fill=orange!12, draw=orange!60!black},
    head/.style  ={block, fill=green!10,  draw=green!50!black},
    io/.style    ={block, fill=gray!10,   draw=black!60, minimum width=5cm},
    arr/.style   ={-{Latex[length=2mm]}, thick},
    tag/.style   ={font=\footnotesize\itshape, text=black!70}
]
% Input
\node[io] (input)  {Input image tensor $\;\;(B,\,3,\,448,\,448)$};

% Backbone block (single node, with internal text)
\node[bb, below=of input, align=center] (bb) {%
    \textbf{ResNet18 backbone (pretrained, ImageNet)}\\[1pt]
    \texttt{conv1, bn1, relu, maxpool}\\
    \texttt{layer1, layer2, layer3, layer4}\\[1pt]
    output: $(B,\,512,\,14,\,14)$ \hfill $\approx 11.2$\,M params};

% Adapter block
\node[adapter, below=of bb, align=center] (ad) {%
    \textbf{Adapter (newly initialised)}\\[1pt]
    \texttt{Conv2d}$(512\!\to\!1024,\,3\!\times\!3,\,s\!=\!2)$ + BN + LeakyReLU\\
    \texttt{Conv2d}$(1024\!\to\!1024,\,3\!\times\!3,\,s\!=\!1)$ + BN + LeakyReLU\\[1pt]
    output: $(B,\,1024,\,7,\,7)$ \hfill $\approx 14.2$\,M params};

% Head block
\node[head, below=of ad, align=center] (hd) {%
    \textbf{Detection head (newly initialised)}\\[1pt]
    \texttt{Flatten} $\to$
    \texttt{Linear}$(50176\!\to\!4096)$ + LeakyReLU + Dropout(0.5)\\
    \texttt{Linear}$(4096\!\to\!4410)$\\[1pt]
    output: $(B,\,4410) = (B,\,S\!\cdot\! S\!\cdot\! (C\!+\!5B))$ \hfill $\approx 223$\,M params};

% Reshape
\node[io, below=of hd] (out) {Per-cell tensor $(B,\,7,\,7,\,90)$};

% Arrows
\draw[arr] (input) -- (bb);
\draw[arr] (bb)    -- (ad);
\draw[arr] (ad)    -- (hd);
\draw[arr] (hd)    -- (out);

% Total params annotation
\node[tag, right=2mm of hd] {\shortstack[l]{Total trainable: \\ $\approx 248$\,M}};
\end{tikzpicture}
\caption{Architecture used in V4--V6. Blue: pretrained ResNet18 trunk loaded
from torchvision (\texttt{ResNet18\_Weights.DEFAULT}). Orange: convolutional
adapter that downsamples spatially from $14\!\times\!14$ to $7\!\times\!7$
and widens channel depth from $512$ to $1024$. Green: the two-layer fully
connected head; the first dense layer (50\,176~$\to$~4\,096) accounts for
roughly $90\,\%$ of the model's trainable parameters. The final tensor is
reshaped to $(B,\,7,\,7,\,90)$ for loss and decoding; the channel layout per
cell is the $C\!=\!80$ COCO layout defined in
Section~\ref{sec:theory:grid}.}
\label{fig:arch}
\end{figure}

\paragraph{Backbone choice.} The original YOLOv1 \citep{redmon2016you} uses a
custom 24-layer Darknet trunk pretrained on ImageNet. Reproducing this
backbone from scratch and then pretraining it on ImageNet was outside the
compute budget of this project. ResNet18 \citep{he2016deep} is the smallest
member of the residual-network family for which torchvision ships a
ready-to-use pretrained checkpoint, and its output stride and channel count
($\times 32$ downsampling, $512$ channels) match a YOLO grid of $S=7$ on a
$448 \times 448$ input. The exchange we accept is a smaller, faster,
already-pretrained trunk in exchange for a less expressive feature backbone
than the paper's; we revisit this trade-off in the Discussion.

\paragraph{Adapter.} ResNet18's \texttt{layer4} produces $(B, 512, 14, 14)$
feature maps for a $448 \times 448$ input. The detection head needs
$(B, 1024, 7, 7)$ to match the original Darknet head's expected size. The
adapter is the cheapest module that bridges the two: a stride-2 convolution
halves the spatial resolution while doubling the channel depth, and a
stride-1 convolution refines the resulting representation. Both convolutions
are followed by batch normalisation~\citep{ioffe2015batch} and a LeakyReLU
activation with negative slope $0.1$, matching the paper's choice in the
detection layers.

\paragraph{Detection head.} The head is identical to the original paper.
Flatten produces a $50\,176$-element vector, the first dense layer maps it
to $4\,096$ with LeakyReLU and dropout~\citep{srivastava2014dropout}, and the
final dense layer produces the $4\,410 = 7 \cdot 7 \cdot 90$ outputs. No
activation is applied to the final layer: the loss operates directly on the
raw logits, with no sigmoid on the box coordinates and no softmax on the
class probabilities.

\paragraph{Parameter accounting.} The first dense layer of the head accounts
for roughly $205$\,M of the $248$\,M trainable parameters
(\,$50\,176 \times 4\,096$\,). The ResNet18 trunk carries about $11.2$\,M, the
adapter contributes $\sim 14$\,M, and the remainder is split between the
final dense layer and the various biases and BN scale/offset terms. The
disparity between the convolutional trunk and the dense head is a known
weakness of the YOLOv1 design and one of the primary motivations for the
fully convolutional rewrite in YOLOv2 and YOLOv3 \citep{redmon2018yolov3}.

% -------------------------------------------------------------
\subsection{Loss implementation}
\label{sec:impl:loss}

The loss is the five-term sum-of-squared-errors formulation from
\citet{redmon2016you}. We implement it as a single \texttt{nn.Module} subclass
(\texttt{src/loss.py}) operating on per-batch tensors. A condensed listing of
the body is given in Listing~\ref{lst:loss-body}; we discuss four
implementation details that are not visible from the mathematical formulation.

\begin{lstlisting}[caption={Body of \texttt{YoloLoss.forward}, abbreviated for
brevity. Auxiliary index variables and the IoU-based responsibility mask are
omitted; the full source is in \texttt{src/loss.py}.},
label=lst:loss-body]
predictions = predictions.reshape(-1, S, S, C + 5*B)

C = self.C
OBJ_IDX    = C
BOX1_SLICE = slice(C + 1, C + 5)
CONF2_IDX  = C + 5
BOX2_SLICE = slice(C + 6, C + 10)

iou_b1   = iou(predictions[..., BOX1_SLICE], target[..., BOX1_SLICE])
iou_b2   = iou(predictions[..., BOX2_SLICE], target[..., BOX1_SLICE])
best_box = torch.max(torch.stack([iou_b1, iou_b2], dim=0), dim=0).indices

exists_box = target[..., OBJ_IDX:OBJ_IDX+1]

# Box loss (xy directly, wh under sign(.)*sqrt(|.|))
box_pred = exists_box * (best_box * predictions[..., BOX2_SLICE]
                       + (1 - best_box) * predictions[..., BOX1_SLICE])
box_tgt  = exists_box * target[..., BOX1_SLICE]

xy_pred       = box_pred[..., 0:2]
wh_pred_sqrt  = torch.sign(box_pred[..., 2:4]) * torch.sqrt(
                    torch.abs(box_pred[..., 2:4]) + 1e-6)
box_pred      = torch.cat([xy_pred, wh_pred_sqrt], dim=-1)
box_tgt       = torch.cat([box_tgt[..., 0:2],
                           torch.sqrt(box_tgt[..., 2:4] + 1e-6)], dim=-1)
box_loss  = mse(box_pred, box_tgt)

# Object + no-object + class losses (see src/loss.py for full code)
\end{lstlisting}

\paragraph{Sum reduction.} We use \texttt{nn.MSELoss(reduction="sum")} rather
than \texttt{"mean"}. The paper specifies a sum of squared errors, and the
relative weighting between the five terms is calibrated for sum reduction:
$\lambda_{\text{coord}} = 5$ is meaningful only if all of the coordinate
summands are accumulated rather than averaged. The consequence is that the
absolute loss magnitude scales with batch size and with grid density; at our
$S = 7$, $B = 2$, $C = 80$ and batch size $32$, the per-batch loss settles in
the low hundreds (as seen in V5 and V6 in
Section~\ref{sec:experiments:setup}).

\paragraph{Safe square-root on width and height.} The paper applies
$\sqrt{w}$ and $\sqrt{h}$ to equalise the gradient contribution of small and
large boxes. The network output is unconstrained, however, so any negative
prediction during early training would crash the square root. We therefore
compute $\operatorname{sign}(x) \cdot \sqrt{|x| + \varepsilon}$ with
$\varepsilon = 10^{-6}$. The sign is preserved so that the gradient through
the responsibility mask points the prediction back into the positive half-line
on the next step, rather than being trapped by an \texttt{abs} symmetry. This
trick is standard in published YOLO reimplementations but the choice of
$\varepsilon$ location matters: $\sqrt{|x| + \varepsilon}$ is safer than
$\sqrt{|x + \varepsilon|}$ near $x = -\varepsilon$.

\paragraph{Avoiding in-place writes on autograd tensors.} An earlier version
of the code transformed \texttt{box\_predictions[..., 2:4]} in place:
\begin{lstlisting}
box_predictions[..., 2:4] = torch.sign(...) * torch.sqrt(...)
\end{lstlisting}
\noindent This silently severs the autograd graph through the modified slice:
no exception is raised, the loss value is correct, but the backward pass
through the $w, h$ pathway returns zero gradients. The current implementation
re-assembles the four-component box tensor via \texttt{torch.cat} along the
channel dimension, which preserves the graph. This was fixed in the same
commit cycle as the bug investigation of
Section~\ref{sec:bug-investigation}; we mention it here because it is part of
the implementation invariants that future modifications must respect.

\paragraph{Indices derived from $C$.} As discussed at length in
Section~\ref{sec:bug:loss-indices}, the indices into the per-cell prediction
tensor are computed from \texttt{self.C} rather than hardcoded. The same
class can therefore be instantiated for any $C$ without code changes, which
both fixes the original COCO bug and keeps the class testable on the
$C\!=\!20$ VOC layout without modifications.

% -------------------------------------------------------------
\subsection{Training infrastructure}
\label{sec:impl:training}

The training loop (\texttt{src/train.py}) is written for multi-GPU operation
under PyTorch's \texttt{DistributedDataParallel} (DDP) and launched with
\texttt{torchrun}. In practice all six runs in this work used a single GPU,
but the multi-process scaffolding (process group init, distributed sampler,
rank-zero-only logging, rank-zero-only checkpointing) was retained so the
loop could be scaled later without restructuring.

\paragraph{Hardware.} V4--V6 ran on a Vast.ai pod with a single NVIDIA H200
NVL (140\,GB HBM3) and node-local NVMe storage. Throughput at $S = 7$,
$B = 2$, $C = 80$, batch size $32$ and input resolution $448\times 448$ was
approximately $114$~seconds per epoch on the training split, or
$\sim\!95$~minutes for $50$ epochs. V1--V3 ran on a different provider
(RunPod) on a pod equipped with eight NVIDIA H100 GPUs, of which only one
was effectively used because the storage substrate became the bottleneck
before compute parallelism could be exploited; the configuration is
discussed under ``Storage substrate, in passing'' below. Absolute timings
for V1--V3 are not directly comparable to V4--V6 because the backbone, the
iteration count, the GPU generation and the I/O path all differ.

\paragraph{Software stack.} PyTorch 2.4.1 with CUDA 12.4 on Python 3.12.
Mixed precision is handled by \texttt{torch.amp.autocast(device\_type="cuda",
dtype=torch.float16)} for the forward pass and the loss, with the optimiser
maintaining a float32 master copy of the weights via
\texttt{torch.amp.GradScaler}. Gradients are unscaled before clipping so that
the clip threshold is interpreted in real magnitude rather than in scaled
units.

\paragraph{Optimiser and schedule.} Stochastic gradient descent with momentum
$0.9$ and weight decay $5\!\times\!10^{-4}$, applied uniformly to every
parameter (we do not use parameter groups with a smaller learning rate on the
backbone). The schedule is a three-epoch linear warm-up from $0$ to the peak
learning rate, followed by cosine annealing to $10^{-5}$
\citep{loshchilov2017sgdr}. The peak rate is $10^{-4}$ in V4--V6 and varies
across V1--V3 as described in Section~\ref{sec:experiments:setup}.

\paragraph{Checkpoints.} The trainer writes three kinds of checkpoint per run:
a rolling \texttt{last.pth} updated every epoch, a snapshot
\texttt{epoch\_NNN.pth} every two epochs, and a \texttt{best.pth} updated only
when the validation loss strictly improves. Each ``full'' checkpoint contains
the model state dict, the optimiser state, the AMP scaler state, the epoch
number and the train/validation losses, so a run can be resumed exactly
(this is how V6 resumes from V5). For evaluation we accept both ``full''
checkpoints and ``clean'' checkpoints that are pure state dicts; the
distinction is detected by the presence of the \texttt{model\_state\_dict} key
and handled in \texttt{scripts/evaluate.py} and
\texttt{scripts/visualize\_prediction.py}.

\paragraph{Storage substrate, in passing.} An earlier attempt to host V1--V3
on a single RunPod pod equipped with $8\!\times$ NVIDIA H100 GPUs and
MooseFS-backed distributed storage was abandoned after the COCO unzip step
repeatedly placed worker processes in uninterruptible (\texttt{D}-state)
sleep, requiring pod reboots and preventing the eight-GPU configuration from
ever being exercised in a meaningful training run. The relevant lesson is
that a $26$~GB dataset of small files benefits substantially from a local
NVMe rather than a network filesystem: switching to a node-local NVMe (which
the Vast.ai pod used for V4--V6 provides natively) was the single largest
engineering improvement of the project, more impactful than any of the
architectural changes made between V1 and V4.

% -------------------------------------------------------------
\subsection{Evaluation pipeline}
\label{sec:impl:eval}

Evaluation is performed offline by \texttt{scripts/evaluate.py} against the
test split. The script loads a checkpoint, iterates over the test
\texttt{DataLoader} with \texttt{shuffle=False} (so that the bookkeeping
indices are deterministic), and for each batch performs the following steps.

\begin{enumerate}
    \item Forward pass under \texttt{torch.no\_grad}, producing the flat
        $(B, 4410)$ output tensor.
    \item Per-cell decoding via \texttt{utils.cellboxes\_to\_boxes}, which
        reshapes to $(B, 7, 7, 90)$, selects the higher-confidence of the two
        candidate boxes per cell, applies the cell-to-image coordinate
        transform discussed in Section~\ref{sec:bug:wh-decoding} (the $1/S$
        factor on $x, y$ only, post-fix), and returns a list of $49$ candidate
        boxes per image in the form
        $[\text{class},\,\text{prob},\,x,\,y,\,w,\,h]$.
    \item Per-class non-maximum suppression, parameterised by a probability
        threshold (the smallest predicted confidence to consider, set to
        $0.1$ and $0.3$ in the reported results) and an IoU threshold (the
        overlap above which a lower-confidence prediction of the same class
        is discarded, fixed at $0.5$).
    \item Aggregation of all surviving predictions across the entire test
        split into a global list of $[\text{image\_idx}, \text{class},
        \text{prob}, x, y, w, h]$ tuples.
    \item Extraction of the corresponding ground-truth list from the target
        tensors, also indexed by \texttt{image\_idx}.
    \item Computation of mean Average Precision via
        \texttt{metrics.mean\_average\_precision}, which sorts each class's
        predictions by confidence, walks them in descending order while
        marking each as a true positive (if it overlaps an unclaimed GT of
        the same class with IoU above the threshold) or false positive, and
        integrates the resulting precision--recall curve per class with the
        trapezoidal rule. The final mAP is the mean of the per-class AP over
        the 80 classes.
\end{enumerate}

\paragraph{One implementation detail.} The mAP code skips classes with zero
ground-truth instances in the evaluated split rather than counting them as
AP~$=$~0. On COCO this never triggers, because every class has at least one
test instance, but it would change the denominator if the same code were
applied to a smaller benchmark.

\paragraph{Reproducibility.} The dataset split is fixed by seed~$42$ in
\texttt{scripts/prepare\_splits.py}, the test data loader uses
\texttt{shuffle=False}, NMS is deterministic, and the loss/decoding code
contains no stochastic elements. Re-running the evaluation script on the
same checkpoint and the same hardware reproduces the reported mAP exactly.
